\documentclass[11pt,a4paper]{article}

\usepackage[a4paper,margin=2.5cm]{geometry}
\usepackage{lmodern}
\usepackage[spanish,es-nodecimaldot]{babel}
\usepackage{microtype}
\usepackage{amsmath,amssymb}
\usepackage{booktabs}
\usepackage{array}
\usepackage{longtable}
\usepackage{tabularx}
\usepackage{enumitem}
\usepackage{hyperref}
\usepackage{xurl}
\usepackage{fancyvrb}
\usepackage{titlesec}

\hypersetup{
  colorlinks=true,
  linkcolor=black,
  urlcolor=blue,
  pdftitle={PyExner - Documentacion de Arquitectura}
}

\setlength{\parindent}{0pt}
\setlength{\parskip}{0.6em}
\setlength{\emergencystretch}{3em}

\titleformat{\section}{\large\bfseries}{\thesection.}{0.5em}{}
\titleformat{\subsection}{\normalsize\bfseries}{\thesubsection.}{0.5em}{}
\titleformat{\subsubsection}{\normalsize\bfseries}{\thesubsubsection.}{0.5em}{}

\newcolumntype{L}[1]{>{\raggedright\arraybackslash}p{#1}}
\newcolumntype{C}[1]{>{\centering\arraybackslash}p{#1}}
\newcolumntype{Y}{>{\raggedright\arraybackslash}X}

\newcommand{\code}[1]{\texttt{#1}}
\newcommand{\pathref}[1]{\href{#1}{\texttt{#1}}}

\begin{document}

\begin{center}
  {\LARGE \textbf{PyExner}}\\[0.4em]
  {\Large Documentacion de Arquitectura}
\end{center}

\section{Vision general}

\textbf{PyExner} es un solver 2D de alto rendimiento para las ecuaciones de aguas someras\newline
(\emph{Shallow Water Equations}, SWE) acopladas con la ecuacion de Exner para transporte de\newline
sedimentos y evolucion morfodinamica del lecho. El codigo esta construido sobre \textbf{JAX} para
aceleracion en GPU y utiliza \textbf{MPI} mediante \code{mpi4py} y \code{mpi4jax} para computacion
distribuida en multiples GPUs o nodos.

\subsection{Ecuaciones que resuelve}

El sistema matematico que PyExner resuelve combina dos fenomenos fisicos: el flujo de agua poco
profunda y la evolucion del lecho por arrastre de sedimento.

\begin{enumerate}[leftmargin=1.8em]
  \item \textbf{Ecuaciones de aguas someras (SWE).}
  Estas ecuaciones describen la conservacion de masa y de momento lineal para un fluido cuya
  profundidad es mucho menor que su extension horizontal. Se obtienen al integrar las ecuaciones
  de Navier-Stokes en la vertical bajo la hipotesis de presion hidrostatica. En su forma
  conservativa bidimensional resultan:

  \begin{align*}
    \frac{\partial h}{\partial t}
    + \frac{\partial (hu)}{\partial x}
    + \frac{\partial (hv)}{\partial y} &= 0 \\
    \frac{\partial (hu)}{\partial t}
    + \frac{\partial}{\partial x}\left(hu^2 + \frac{1}{2}gh^2\right)
    + \frac{\partial (huv)}{\partial y} &= -gh\frac{\partial z}{\partial x} - ghS_{f,x} \\
    \frac{\partial (hv)}{\partial t}
    + \frac{\partial (huv)}{\partial x}
    + \frac{\partial}{\partial y}\left(hv^2 + \frac{1}{2}gh^2\right) &= -gh\frac{\partial z}{\partial y} - ghS_{f,y}
  \end{align*}

  donde $h$ es la profundidad del agua, $u$ y $v$ las componentes de velocidad,
  $g$ la aceleracion de la gravedad, $z$ la cota del lecho y $S_{f}$ la pendiente de friccion
  (modelada con la ley de Manning).

  \item \textbf{Ecuacion de Exner.}
  Describe la evolucion temporal de la cota del lecho erodible $z_b$ como resultado del
  transporte de sedimento de fondo $\mathbf{q_b}$. El factor $1/(1-p)$ corrige por la
  porosidad $p$ de la mezcla granular:

  \[
    \frac{\partial z_b}{\partial t} + \frac{1}{1-p}\nabla \cdot \mathbf{q_b} = 0
  \]
\end{enumerate}

\subsection{Stack tecnologico}

La eleccion de JAX como motor de computo permite compilar funciones numericas a codigo
nativo de GPU (o CPU) mediante \code{jax.jit}, lo que evita la sobrecarga del interprete
de Python en cada iteracion del bucle temporal. La compatibilidad de JAX con
diferenciacion automatica abre, ademas, la puerta a futuros estudios de sensibilidad o
optimizacion adjunta. Para simulaciones de gran escala, el dominio se descompone
espacialmente entre multiples procesos MPI, cada uno con acceso a su propia GPU; la biblioteca
\code{mpi4jax} integra las comunicaciones MPI directamente dentro de las funciones
JIT-compiladas.

\begin{tabularx}{\textwidth}{L{4.2cm}Y}
\toprule
Componente & Tecnologia \\
\midrule
Computo numerico & JAX (\code{jax}, \code{jax.numpy}) \\
Paralelismo distribuido & MPI (\code{mpi4py}, \code{mpi4jax}) \\
I/O paralelo & PnetCDF (\code{pnetcdf}) \\
Configuracion & YAML (\code{pyyaml}) \\
Visualizacion & Matplotlib, xarray \\
\bottomrule
\end{tabularx}

\section{Punto de entrada: \code{run\_driver()}}

\textbf{Archivo:} \code{src/PyExner/runtime/driver.py}

La funcion principal orquesta toda la simulacion. Actua como punto unico de entrada: recibe la ruta
al archivo de configuracion YAML, inicializa todos los subsistemas en el orden correcto y lanza el
bucle temporal. Su flujo logico es el siguiente:

\begin{Verbatim}[fontsize=\small]
config YAML -> Parallel (MPI) -> PnetCDF I/O -> Mesh -> State -> Boundaries -> Solver 
-> Integrator -> Ejecucion
\end{Verbatim}

\subsection{Pasos detallados}

\begin{enumerate}[leftmargin=1.8em]
  \item Lee el archivo YAML de configuracion mediante \code{config\_path}.
  \item Inicializa MPI mediante \code{Parallel(params)} y crea la topologia cartesiana.
  \item Crea el lector PnetCDF para abrir archivos de entrada y salida en paralelo.
  \item Genera la malla con \code{pnetcdf\_io.generate\_mesh()}.
  \item Crea el estado vacio segun el esquema de flujo, \code{Roe} o \code{Roe Exner}.
  \item Lee las condiciones iniciales desde el archivo NetCDF hacia el estado.
  \item Configura las condiciones de contorno mediante \code{BoundaryManager}.
  \item Selecciona el solver desde \code{SOLVER\_REGISTRY}.
  \item Selecciona el integrador temporal desde \code{INTEGRATOR\_REGISTRY}.
  \item Ejecuta el bucle temporal mediante \code{integrator.run\_fn()}.
\end{enumerate}

\subsection{Parametros YAML principales}

\begin{longtable}{L{3.2cm}L{8.2cm}L{3cm}}
\toprule
Parametro & Descripcion & Ejemplo \\
\midrule
\endhead
\code{end\_time} & Tiempo final de simulacion & \code{20.0} \\
\code{out\_freq} & Frecuencia de escritura de salida & \code{0.25} \\
\code{cfl} & Numero de Courant-Friedrichs-Lewy. Controla el paso de tiempo para garantizar estabilidad numerica: valores cercanos a 1 son mas eficientes pero menos robustos & \code{0.5} \\
\code{flux\_scheme} & Esquema numerico: \code{Roe} o \code{Roe Exner} & \code{Roe Exner} \\
\code{integrator} & Esquema temporal: \code{Forward Euler} o \code{SSPRK2} & \code{Forward Euler} \\
\code{parNx}, \code{parNy} & Particiones MPI en X e Y & \code{2}, \code{2} \\
\code{input\_file} & Archivo NetCDF de entrada & \code{input.nc} \\
\code{output\_file} & Archivo NetCDF de salida & \code{output.nc} \\
\code{boundaries} & Definicion de condiciones de contorno & ver Seccion 6 \\
\code{erosion} & Parametros de sedimentos para \code{Roe Exner} & ver Seccion 4 \\
\bottomrule
\end{longtable}

\section{Modulo \code{state/}: variables de estado}

\textbf{Directorio:} \code{src/PyExner/state/}

El modulo \code{state/} define las variables conservativas que el solver transporta en cada celda
de la malla. Se denominan \emph{conservativas} porque representan cantidades cuyo integral se
conserva globalmente (masa de agua $h$, momento $hu$, $hv$), a diferencia de las variables
\emph{primitivas} ($h$, $u$, $v$) que no satisfacen directamente las leyes de conservacion
en forma de divergencia.

\subsection{\code{BaseState}}

\textbf{Archivo:} \code{src/PyExner/state/base.py}

Clase base con los campos hidraulicos comunes:

\begin{tabularx}{\textwidth}{L{3cm}L{3cm}Y}
\toprule
Campo & Tipo & Descripcion \\
\midrule
\code{h} & \code{jax.Array} & Profundidad del agua \\
\code{hu} & \code{jax.Array} & Momento en X, es decir, $h \cdot u$ \\
\code{hv} & \code{jax.Array} & Momento en Y, es decir, $h \cdot v$ \\
\bottomrule
\end{tabularx}

Metodos utilitarios relevantes: \code{unshard()}, \code{\_\_getitem\_\_()}, \code{reshape()},
\code{apply\_to\_all()} y \code{to\_host()}.

\subsection{\code{RoeState}}

\textbf{Archivo:} \code{src/PyExner/state/roe\_state.py}

Este estado se utiliza para simulaciones solo hidraulicas.

\begin{tabularx}{\textwidth}{L{4cm}Y}
\toprule
Campo adicional & Descripcion \\
\midrule
\code{z} & Elevacion del lecho fija (no cambia durante la simulacion) \\
\code{n} & Coeficiente de rugosidad de Manning \\
\bottomrule
\end{tabularx}

El coeficiente de Manning $n$ parametriza la resistencia al flujo que ejerce el lecho. Valores
tipicos van de $\sim$0.01 para canales lisos a $\sim$0.05 para cauces naturales con vegetacion.

Metodos principales: \code{empty(mesh)}, \code{from\_params(params)} y \code{replace(**kwargs)}.

\subsection{\code{RoeExnerState}}

\textbf{Archivo:} \code{src/PyExner/state/roe\_exner\_state.py}

Este estado se utiliza para simulaciones hidrodinamicas y morfodinamicas acopladas. Distingue
entre una capa de roca fija (\code{z}) y una capa de sedimento suelto (\code{z\_b}) que puede
erosionarse o depositarse. La superficie del lecho que percibe el flujo es la suma de ambas.

\begin{longtable}{L{3.4cm}L{11cm}}
\toprule
Campo adicional & Descripcion \\
\midrule
\endhead
\code{z} & Elevacion del sustrato no erodible (roca base) \\
\code{z\_b} & Elevacion del lecho erodible (capa de sedimento) \\
\code{n} & Rugosidad de Manning original \\
\code{n\_b} & Rugosidad efectiva calculada a partir de $d_{50}$ \\
\code{G} & Factor de interaccion morfodinamica para la tasa de transporte \\
\code{seds} & Propiedades de sedimentos: fraccion, diametro, densidad y afines \\
\bottomrule
\end{longtable}

\subsection{Registro de estados}

\textbf{Archivo:} \code{src/PyExner/state/registry.py}

El sistema de registro usa decoradores:

\begin{itemize}[leftmargin=1.6em]
  \item \code{@register\_state(Roe)} registra \code{RoeState}.
  \item \code{@register\_state(Roe Exner)} registra \code{RoeExnerState}.
  \item \code{create\_empty\_state(name, mesh, rank)} crea un estado vacio por nombre.
  \item \code{create\_state(name, params)} crea una instancia a partir de parametros.
\end{itemize}

\section{Modulo \code{solvers/}: solvers de Riemann}

\textbf{Directorio:} \code{src/PyExner/solvers/}

Un solver de Riemann calcula el flujo numerico en la interfaz entre dos celdas a partir de los
estados a cada lado. Resuelve (de forma exacta o aproximada) el problema de Riemann, es decir,
la evolucion de una discontinuidad inicial entre dos estados constantes. El resultado determina
cuanto de cada variable conservativa pasa de una celda a la otra en un paso temporal.

\subsection{Arquitectura del registro}

\textbf{Archivo:} \code{src/PyExner/solvers/registry.py}

Cada solver se registra como un \code{SolverBundle} con los siguientes campos:

\begin{tabularx}{\textwidth}{L{3cm}L{3cm}Y}
\toprule
Campo & Tipo & Descripcion \\
\midrule
\code{name} & \code{str} & Nombre del solver \\
\code{config} & \code{Callable} & Funcion de configuracion que crea \code{SolverConfig} \\
\code{mask\_fn} & \code{Callable} & Genera mascaras de celdas activas o inactivas \\
\code{init\_fn} & \code{Callable} & Inicializa el estado e intercambia halos \\
\code{step\_fn} & \code{Callable} & Avanza un paso temporal \\
\code{compute\_dt\_fn} & \code{Callable} & Calcula $\Delta t$ segun la condicion CFL \\
\bottomrule
\end{tabularx}

\code{SolverConfig} contiene \code{mpi\_handler}, \code{boundaries}, \code{dx} y
\code{halo\_exchange}.

\subsection{Solver de Roe para SWE}

\textbf{Archivo del solver:} \code{src/PyExner/solvers/roe\_solver.py}\\
\textbf{Kernels numericos:} \code{src/PyExner/solvers/kernels/roe.py}

Este solver resuelve las ecuaciones de aguas someras 2D sin evolucion del lecho.

\subsubsection{Funcion principal: \code{roe\_solver(si, sj, nx, ny, dx)}}

La funcion implementa el solver de Riemann aproximado de Roe en la interfaz entre dos celdas vecinas
\code{si} y \code{sj}. El procedimiento puede resumirse asi:

\begin{enumerate}[leftmargin=1.8em]
  \item Rota las velocidades al marco normal-tangencial de la interfaz. Esto permite tratar
  la interfaz como un problema unidimensional independiente de la orientacion de la malla.
  \item Calcula los promedios de Roe. Son valores intermedios especialmente elegidos para que la
  linializacion del flujo sea consistente con las relaciones de salto de Rankine-Hugoniot:
  \[
    \tilde{h} = \frac{1}{2}(h_i + h_j), \qquad
    \tilde{u} = \frac{\sqrt{h_i}\hat{u}_i + \sqrt{h_j}\hat{u}_j}{\sqrt{h_i} + \sqrt{h_j}}, \qquad
    \tilde{c} = \sqrt{g\tilde{h}}
  \]
  \item Obtiene los autovalores, que representan las velocidades de propagacion de las tres
  familias de ondas del sistema (onda lenta, contacto y onda rapida):
  \[
    \lambda_1 = \tilde{u} - \tilde{c}, \qquad
    \lambda_2 = \tilde{u}, \qquad
    \lambda_3 = \tilde{u} + \tilde{c}
  \]
  \item Aplica la correccion de entropia de Harten-Hyman. Esta correccion evita que el esquema
  produzca ondas de rarefaccion no fisicas (choques expansivos) cuando un autovalor se acerca
  a cero; sustituye el valor absoluto $|\lambda|$ por una funcion suavizada en un entorno
  del origen.
  \item Construye los autovectores y la descomposicion modal. La matriz $P$ de autovectores
  derechos permite descomponer el salto de estado $\Delta U$ en contribuciones independientes
  de cada familia de ondas.
  \item Calcula las ondas mediante $\boldsymbol{\alpha} = P^{-1}\Delta U$.
  \item Incorpora terminos fuente \emph{well-balanced}. Un esquema se llama
  \emph{well-balanced} cuando es capaz de preservar estados estacionarios exactos como
  el agua en reposo sobre batimetria irregular. Sin este tratamiento, los errores de
  truncacion en el termino de pendiente del lecho generan corrientes espurias.
    \begin{itemize}
      \item \textbf{Thrust topografico:} formulacion especial para mantener el equilibrio en
            reposo sobre batimetria variable, basada en Murillo y Garc\'{\i}a-Navarro (2010).
      \item \textbf{Friccion de Manning:} $\tau = g\tilde{h}\,S_f\,\Delta x$ con
            $S_f = \dfrac{n^2\sqrt{u^2+v^2}\cdot u}{h^{4/3}}$.
    \end{itemize}
  \item Aplica \emph{upwinding} para separar contribuciones positivas y negativas.
  \item Rota de regreso al marco global $(x, y)$.
\end{enumerate}

\subsubsection{Paso bidimensional completo}

La funcion \code{roe\_solve\_2D(state, dt, dx)} crea los cortes de interfaces en X e Y, llama a
\code{roe\_solver()} en ambas direcciones y actualiza el estado segun

\[
  U^{n+1} = U^n - \frac{\Delta t}{\Delta x}F
\]

\subsubsection{Calculo del paso de tiempo}

Para garantizar la estabilidad numerica, el paso de tiempo se adapta en cada iteracion segun la
condicion de Courant-Friedrichs-Lewy (CFL). Esta condicion exige que ninguna onda fisicamente
relevante recorra mas de una celda por paso de tiempo. La funcion
\code{compute\_dt(state, mask, dx)} lo calcula como:

\[
  \Delta t = \mathrm{CFL} \times \min\left(\frac{\Delta x}{|u| + c},\frac{\Delta y}{|v| + c}\right)
\]

El valor se reduce globalmente con \code{MPI.MIN}, de modo que todos los procesos usan el
paso de tiempo mas restrictivo del dominio completo.

\subsubsection{Flujo del \code{step\_fn\_roe}}

\begin{Verbatim}[fontsize=\small]
Estado -> roe_solve_2D -> Intercambio de halos (MPI)
-> Aplicar contorno -> Estado actualizado
\end{Verbatim}

\subsection{Solver Roe-Exner para SWE + Exner}

\textbf{Archivo del solver:} \code{src/PyExner/solvers/roe\_exner\_solver.py}\\
\textbf{Kernels numericos:} \code{src/PyExner/solvers/kernels/roe\_exner.py}

Este solver resuelve las ecuaciones de aguas someras acopladas con la ecuacion de Exner.

\subsubsection{Diferencias clave respecto a Roe}

\begin{enumerate}[leftmargin=1.8em]
  \item El estado incluye \code{z\_b}, \code{G}, \code{n\_b} y \code{seds}.
  \item La funcion \code{compute\_G()} estima la tasa de transporte con una formulacion tipo
  Meyer-Peter y Muller (MPM) modificada, con transporte $q_b \propto \gamma_1 \cdot \gamma_2 \cdot \gamma_3$.
  Soporta multiples fracciones de sedimento con mezcla ponderada.
  \item \code{exner\_solver} implementa el \emph{Approximately Coupled Method} (ACM), basado en
  Mart\'{\i}nez-Aranda et al.\ (2021), que calcula el flujo de sedimento en interfaces usando
  \emph{upwinding} basado en $\lambda_4$ (velocidad de onda de Exner) y actualiza la elevacion del lecho:
  \[
    z_b^{n+1} = z_b^n - \frac{\Delta t}{\Delta x}\nabla \cdot \mathbf{q_b}
  \]
  \item \code{\_get\_lambda} resuelve el polinomio cubico del sistema acoplado SWE-Exner
  mediante las formulas trigonometricas de Cardano-Vieta para obtener $\lambda_{\min}$ y $\lambda_{\max}$.
  \item \code{momentum\_corrections} anula el momento en direcciones donde el agua no puede fluir
  (celdas secas adyacentes con lecho mas alto que la superficie libre, o paredes).
\end{enumerate}

\subsubsection{Expresiones relevantes para \code{G}}

El factor \code{G} mide la intensidad del transporte de sedimento de fondo. Se calcula a partir
del parametro de Shields $\theta$, que adimensionaliza la tension tangencial del flujo sobre una
particula de sedimento. Si $\theta$ supera un umbral critico $\theta_c$, el sedimento se pone
en movimiento:

\[
  \theta = \frac{n_p^2(u^2 + v^2)}{(s-1)d_{50}h^{1/3}},
  \qquad
  \Delta \theta = \max(\theta - \theta_c, 0)
\]

donde $n_p$ es la rugosidad de Manning asociada a la particula, $s = \rho_s / \rho_w$ es la
densidad relativa del sedimento y $d_{50}$ el diametro mediano del grano. El exceso
$\Delta\theta$ alimenta la formula de transporte tipo MPM, que estima el caudal solido por
unidad de ancho. Cuando existen multiples fracciones granulometricas, el transporte total se
obtiene como una suma ponderada por la fraccion en masa de cada clase.

\subsubsection{Flujo del \code{step\_fn\_roeexner}}

\begin{Verbatim}[fontsize=\small]
1. roe_solve_2D         -> Resolver hidrodinamica (SWE)
2. momentum_corrections -> Corregir momento en celdas bloqueadas
3. boundaries.apply     -> Aplicar condiciones de contorno
4. halo_exchange        -> Sincronizar h, hu, hv entre procesos MPI
5. compute_G + halo     -> Calcular y sincronizar el factor G
6. exner_solve_2D       -> Resolver evolucion del lecho
7. halo_exchange z_b    -> Sincronizar lecho y recalcular n_b
\end{Verbatim}

\subsection{HLLC reservado}

\textbf{Archivo:} \code{src/PyExner/solvers/kernels/hllc.py}

El archivo existe como reserva para una implementacion futura del solver HLLC.

\section{Modulo \code{integrators/}: integradores temporales}

\textbf{Directorio:} \code{src/PyExner/integrators/}

\subsection{Registro}

\textbf{Archivo:} \code{src/PyExner/integrators/registry.py}

Cada integrador se representa como un \code{IntegratorBundle}:

\begin{tabularx}{\textwidth}{L{3cm}L{3cm}Y}
\toprule
Campo & Tipo & Descripcion \\
\midrule
\code{name} & \code{str} & Nombre del integrador \\
\code{config} & \code{Callable} & Construye \code{IntegratorConfig} \\
\code{run\_fn} & \code{Callable} & Bucle temporal principal \\
\bottomrule
\end{tabularx}

\code{IntegratorConfig} contiene \code{cfl}, \code{end\_time}, \code{out\_freq},
\code{solver\_config} y \code{solver\_bundle}.

\subsection{Forward Euler}

\textbf{Archivo:} \code{src/PyExner/integrators/forwardeuler.py}

Es el esquema temporal mas simple: de primer orden en tiempo con error de truncacion
$\mathcal{O}(\Delta t)$. Resulta util para prototipado rapido o problemas donde la precision
temporal no es el factor limitante. Su estructura interna combina un bucle Python convencional
(para la escritura de salida) con un bucle JIT-compilado por JAX (para las iteraciones
numericas, sin sobrecarga de interprete):

\begin{enumerate}[leftmargin=1.8em]
  \item Un bucle exterior en Python controla escritura de salida y tiempos de salida.
  \item Un bucle interior con \code{jax.lax.while\_loop} ejecuta iteraciones JIT hasta el siguiente tiempo de salida.
  \item Un paso final ajusta $\Delta t$ para caer exactamente sobre el tiempo de salida.
  \item El estado se escribe con \code{io.write\_state()} en formato PnetCDF paralelo.
\end{enumerate}

\subsection{SSPRK2}

\textbf{Archivo:} \code{src/PyExner/integrators/SSPRK2.py}

El esquema \emph{Strong Stability Preserving} Runge-Kutta de segundo orden ofrece doble orden
de precision temporal respecto a Forward Euler manteniendo la propiedad TVD
(\emph{Total Variation Diminishing}): si el operador espacial no aumenta la variacion total
con un paso de Euler, tampoco lo hara el paso completo SSPRK2. Esto es especialmente importante
en problemas con choques (roturas de presa, frentes de onda) donde preservar la monotonicidad
evita oscilaciones espurias.

El esquema se expresa como:

\begin{align*}
  U^{(1)} &= U^n + \Delta t\,L(U^n) \\
  U^{n+1} &= \frac{1}{2}U^n + \frac{1}{2}U^{(1)} + \frac{1}{2}\Delta t\,L(U^{(1)})
\end{align*}

Es equivalente a un promedio entre el estado original y dos pasos tipo Euler. La propiedad SSP
garantiza que el paso completo es estable bajo la misma restriccion CFL que un solo paso de Euler.

\section{Modulo \code{domain/}: malla y condiciones de contorno}

\textbf{Directorio:} \code{src/PyExner/domain/}

\subsection{Malla}

\textbf{Archivo:} \code{src/PyExner/domain/mesh.py}

\code{Mesh2D} es un \emph{dataclass} para una malla estructurada uniforme. Se asume que el
espaciamiento es el mismo en ambas direcciones ($\Delta x = \Delta y$), lo que simplifica el
calculo de flujos y evita factores de correccion geometrica en el solver de Riemann.

\begin{tabularx}{\textwidth}{L{4cm}Y}
\toprule
Campo & Descripcion \\
\midrule
\code{global\_Ny}, \code{global\_Nx} & Dimensiones globales del dominio \\
\code{local\_Ny}, \code{local\_Nx} & Dimensiones locales por proceso MPI \\
\code{x\_offset}, \code{y\_offset} & Desplazamiento del subdominio local \\
\code{local\_X}, \code{local\_Y} & Coordenadas locales \\
\code{dh} & Espaciamiento de malla con $\Delta x = \Delta y$ \\
\code{local\_shape}, \code{global\_shape} & Tuplas con las formas local y global \\
\bottomrule
\end{tabularx}

La malla se genera automaticamente a partir del archivo NetCDF de entrada y se divide entre procesos MPI.

\subsection{Condiciones de contorno}

\textbf{Archivo:} \code{src/PyExner/domain/boundary\_registry.py}

\code{BoundaryManager} es la clase encargada de gestionar las condiciones de contorno. La
identificacion de celdas fronterizas se realiza geometricamente: el usuario define poligonos
en el archivo YAML y el codigo comprueba, mediante un test de punto-en-poligono compilado
con JIT, que celdas caen dentro de cada zona de contorno. Este enfoque es mas flexible que
imponer condiciones sobre filas o columnas fijas, pues permite fronteras arbitrarias (curvas,
obstaculos internos, etc.).

Sus responsabilidades principales son:

\begin{itemize}[leftmargin=1.6em]
  \item Leer las definiciones de contorno desde YAML.
  \item Identificar celdas de contorno mediante los poligonos mencionados.
  \item Calcular indices interiores adyacentes para condiciones reflectivas o transmisivas.
  \item Aplicar todas las condiciones de contorno secuencialmente con \code{apply(state, time)}.
\end{itemize}

\subsection{Tipos de condicion de contorno}

\textbf{Directorio:} \code{src/PyExner/domain/boundaries/}

{\footnotesize
\begin{longtable}{L{5cm}L{4.5cm}L{5.2cm}}
\toprule
Clase & Registro & Descripcion \\
\midrule
\endhead
\code{RoeReflectiveBoundary} & \code{Roe Reflective} & Pared solida: copia escalares y refleja la componente normal del momento \\
\code{RoeTransmissiveBoundary} & \code{Roe Transmissive} & Salida libre: copia todos los campos desde la celda interior \\
\code{RoeExnerReflectiveBoundary} & \code{Roe Exner Reflective} & Pared solida con campos adicionales de Exner \\
\code{RoeExnerTransmissiveBoundary} & \code{Roe Exner Transmissive} & Salida libre con campos de Exner \\
\code{RoeExnerZeroMomentumBoundary} & \code{Roe Exner ZeroMomentum} & Impone \code{hu = hv = 0} en la frontera \\
\bottomrule
\end{longtable}
}

Todas las clases de contorno se registran como \emph{pytree} de JAX para ser compatibles con
\code{jax.jit}. A continuacion se describe brevemente cada tipo:

\begin{itemize}[leftmargin=1.6em]
  \item \textbf{Reflectiva:} modela una pared solida impermiable. Se copian los escalares de la
  celda interior adyacente pero se invierte el signo de la componente normal del momento, de modo
  que el flujo normal neto es cero.
  \item \textbf{Transmisiva:} modela una salida libre. Todos los campos se copian directamente
  desde la celda interior, permitiendo que las ondas abandonen el dominio sin generar reflexiones
  artificiales.
  \item \textbf{Zero Momentum} (solo Roe Exner): impone $hu = hv = 0$ en las celdas fronterizas,
  creando una zona de remanso sin velocidad.
\end{itemize}

\subsection{Ejemplo de definicion YAML}

\begin{Verbatim}[fontsize=\small]
boundaries:
  wall_left:
    type: Reflective
    polygon:
      - [0.0, 0.0]
      - [0.01, 0.0]
      - [0.01, 10.0]
      - [0.0, 10.0]
    values: [NaN, 0.0, NaN, NaN]
    normal: [-1.0, 0.0]
\end{Verbatim}

El campo \code{polygon} define los vertices de la zona fronteriza en coordenadas del dominio.
El vector \code{values} permite forzar valores concretos en las celdas de contorno; los campos
marcados como \code{NaN} no se modifican y se tratan segun la logica propia del tipo de condicion.
El vector \code{normal} indica la orientacion de la pared (normal exterior unitaria), necesario
para la reflexion correcta del momento.

El valor de \code{flux\_scheme} se concatena con \code{type} para consultar el registro. Por ejemplo,
\code{Roe Exner} mas \code{Reflective} produce \code{Roe Exner Reflective}.

\section{Modulo \code{parallel/}: computacion distribuida MPI}

\textbf{Archivo:} \code{src/PyExner/parallel/mpi\_utils.py}

Cuando el dominio es demasiado grande para caber en la memoria de una sola GPU, o cuando se
desea reducir el tiempo de computo, PyExner descompone la malla en subdominios rectangulares y
asigna uno a cada proceso MPI. Los procesos se organizan en una topologia cartesiana 2D que
refleja la disposicion espacial del dominio, facilitando la identificacion de vecinos.

\subsection{Clase \code{Parallel}}

\begin{tabularx}{\textwidth}{L{3.5cm}Y}
\toprule
Atributo & Descripcion \\
\midrule
\code{parNx}, \code{parNy} & Numero de particiones en X e Y \\
\code{cart\_comm} & Comunicador cartesiano 2D \\
\code{rank}, \code{size} & Rango del proceso y numero total de procesos \\
\code{coords} & Coordenadas del proceso en la topologia \\
\code{dims} & Dimensiones de la topologia, $(parNy, parNx)$ \\
\code{neighbors} & Rangos vecinos: \code{north}, \code{south}, \code{east}, \code{west} \\
\bottomrule
\end{tabularx}

\subsection{Intercambio de halos}

Cada subdominio necesita conocer una franja de celdas pertenecientes a sus vecinos para evaluar
correctamente los flujos en las interfaces compartidas. Estas celdas auxiliares se denominan
\emph{halos} o \emph{ghost cells}. Tras cada actualizacion del estado, los halos se sincronizan
mediante envios punto a punto de MPI.

El intercambio de halos se define en cada kernel de solver, por ejemplo en
\code{src/PyExner/solvers/\allowbreak kernels/roe.py} y
\code{src/PyExner/solvers/\allowbreak kernels/roe\_exner.py}.

\begin{itemize}[leftmargin=1.6em]
  \item Usa \code{mpi4jax} para llamadas MPI compatibles con JAX y JIT, de modo que las
  comunicaciones se integran en el grafo de computo sin salir de la compilacion JIT.
  \item Intercambia celdas fantasma de ancho 1 (suficiente para un stencil de 3 puntos).
  \item El orden de envio es: \code{west -> north -> east -> south}.
  \item El orden de recepcion es: \code{east -> south -> west -> north}.
  \item El orden circular evita bloqueos mutuos (\emph{deadlock}) al garantizar que cada
  proceso envia antes de intentar recibir del mismo vecino.
\end{itemize}

\section{Modulo \code{io/}: entrada y salida}

\textbf{Directorio:} \code{src/PyExner/io/}

PyExner utiliza PnetCDF (Parallel NetCDF) como formato de archivo. A diferencia del NetCDF
clasico, PnetCDF permite que multiples procesos MPI lean y escriban simultaneamente el mismo
archivo, evitando cuellos de botella en la serializacion de datos y escalando correctamente a
cientos de procesos.

\subsection{\code{PnetCDFStateIO}}

\textbf{Archivo:} \code{src/PyExner/io/pnetcdf\_reader.py}

Esta clase centraliza la lectura y escritura en formato PnetCDF.

\begin{tabularx}{\textwidth}{L{4.5cm}Y}
\toprule
Metodo & Descripcion \\
\midrule
\code{generate\_mesh()} & Lee coordenadas globales, divide el dominio y retorna un \code{Mesh2D} \\
\code{read\_state(state, mesh, config)} & Lee variables del NetCDF de entrada y llena el estado, incluyendo halos si aplica \\
\code{write\_state(state, mesh, mask)} & Escribe el estado actual al NetCDF de salida de forma colectiva \\
\bottomrule
\end{tabularx}

Tambien soporta lectura de parametros sedimentologicos desde YAML, como fraccion, diametro $d_{50}$,
densidad, flujos de erosion y deposicion, y porosidad de la mezcla.

\subsection{Archivos reservados}

\begin{itemize}[leftmargin=1.6em]
  \item \code{src/PyExner/io/diagnostics.py}: reservado para diagnosticos futuros.
  \item \code{src/PyExner/io/visualizer.py}: reservado para visualizacion futura.
\end{itemize}

\section{Modulo \code{utils/}: constantes y utilidades}

\textbf{Directorio:} \code{src/PyExner/utils/}

\subsection{Constantes}

\textbf{Archivo:} \code{src/PyExner/utils/constants.py}

El modulo define tolerancias numericas que evitan divisiones por cero y filtran valores
fisicamente irrelevantes. \code{DRY\_TOL} marca el limite por debajo del cual una celda se
considera seca (sin flujo); \code{SED\_TOL} previene el transporte de cantidades infinitesimales
de sedimento; \code{VEL\_TOL} anula velocidades residuales menores que el ruido numerico;
y \code{TIMESTEP\_TOL} y \code{FLUX\_TOL} proporcionan margenes de seguridad para el
integrador y los flujos interfaz, respectivamente.

\begin{tabularx}{\textwidth}{L{3.2cm}L{2.2cm}Y}
\toprule
Constante & Valor & Descripcion \\
\midrule
\code{g} & \code{9.81} & Aceleracion de la gravedad \\
\code{DRY\_TOL} & \code{1e-3} & Tolerancia para celdas secas \\
\code{SED\_TOL} & \code{1e-4} & Tolerancia para sedimento minimo \\
\code{VEL\_TOL} & \code{1e-6} & Tolerancia para velocidades pequenas \\
\code{TIMESTEP\_TOL} & \code{5e-7} & Tolerancia para el integrador temporal \\
\code{FLUX\_TOL} & \code{1e-12} & Tolerancia para flujos \\
\bottomrule
\end{tabularx}

\section{Tests y casos de ejemplo}

\textbf{Directorio:} \code{tests/}

Los casos de prueba cubren problemas clasicos de validacion en hidraulica computacional.
Las roturas de presa (\emph{dam break}) son referencias estandar porque admiten solucion
analitica o semiexacta, lo que permite verificar la convergencia y la correcta captura de
choques. Los casos con sedimento (\code{Roe Exner}) validan la evolucion del lecho comparando
con resultados publicados.

{\footnotesize
\begin{longtable}{L{2.8cm}L{6.1cm}L{1.8cm}L{3.5cm}}
\toprule
Caso & Directorio & Esquema & Descripcion \\
\midrule
\endhead
\textbf{1D Dam Break} & \code{tests/runtime/1D\_dambreak/} & Roe & Rotura de presa 1D clasica \\
\textbf{2D Dam Break} & \code{tests/runtime/2D\_dambreak/} & Roe & Rotura de presa 2D simetrica \\
\textbf{Symmetrical Dam Break} & \code{tests/runtime/symmetrical\_dambreak/} & Roe & Caso de dam break simetrico \\
\textbf{Erodible Channel} & \code{tests/runtime/erodible\_channel/} & Roe Exner & Canal erodible con transporte de sedimentos \\
\textbf{L-Domain} & \code{tests/runtime/L\_domain/} & Roe Exner & Dominio en forma de L con erosion \\
\bottomrule
\end{longtable}
}

Cada caso incluye, en general, un script para construir el input, un archivo YAML de configuracion,
un script de ejecucion y un script de analisis de datos.

\section{Diagrama de flujo de una simulacion}

\begin{Verbatim}[fontsize=\small]
run_driver()
|
+- 1. Leer YAML config
+- 2. Parallel(params) <- Topologia MPI cartesiana
+- 3. PnetCDFStateIO <- Abrir archivos NC
+- 4. generate_mesh() <- Dividir dominio
+- 5. create_empty_state() <- RoeState o RoeExnerState
+- 6. read_state() <- Leer condiciones iniciales
+- 7. BoundaryManager() <- Configurar contornos
+- 8. create_solver_bundle() <- Roe o Roe Exner
+- 9. create_integrator_bundle() <- FE o SSPRK2
|
+- 10. integrator.run_fn()
    |
    +- while time < end_time:
       +- jax.lax.while_loop:
          +- compute_dt (CFL)
          +- solver.step_fn()
             +- Roe solve 2D (SWE)
             +- [Momentum corrections]
             +- Boundary apply
             +- Halo exchange (MPI)
             +- [compute_G + Exner]
             +- [Halo z_b + n_b]
          +- time += dt
       +- io.write_state() <- salida NC

return state, (X, Y)
\end{Verbatim}

Los pasos entre corchetes solo aplican al solver \code{Roe Exner}.

\section{Sistema de registros}

PyExner usa un \emph{patron de registro con decoradores} para facilitar extensibilidad. En lugar
de que el codigo central contenga sentencias \code{if/elif} que enumeren todas las opciones, cada
nuevo componente se registra automaticamente al importarse. Esto desacopla la logica de seleccion
de la implementacion concreta y permite agregar variantes sin modificar el codigo existente
(principio abierto/cerrado). Los registros principales son:

{\footnotesize
\begin{tabularx}{\textwidth}{L{3.2cm}L{6cm}Y}
\toprule
Registro & Decorador & Uso \\
\midrule
\code{STATE\_REGISTRY} & \code{@register\_state(nombre)} & Estados o variables conservativas \\
\code{SOLVER\_REGISTRY} & \code{@register\_solver\_bundle(nombre)} & Solvers de Riemann \\
\code{INTEGRATOR\_REGISTRY} & \code{@register\_integrator\_bundle(nombre)} & Integradores temporales \\
\code{BOUNDARY\_REGISTRY} & \code{@register\_boundary(nombre)} & Condiciones de contorno \\
\bottomrule
\end{tabularx}
}

Para agregar un nuevo estado, solver, integrador o condicion de contorno, basta con crear la nueva
clase o funcion, decorarla adecuadamente e importarla desde el \code{\_\_init\_\_.py} correspondiente.

\end{document}